% IEEE Conference template — IS 584 Final Report
% Tip-of-the-Tongue Retrieval on TREC 2023
\documentclass[conference]{IEEEtran}
\IEEEoverridecommandlockouts
\usepackage{cite}
\usepackage{amsmath,amssymb,amsfonts}
\usepackage{algorithmic}
\usepackage{graphicx}
\usepackage{textcomp}
\usepackage{xcolor}
\usepackage{booktabs}
\usepackage{url}
\def\BibTeX{{\rm B\kern-.05em{\sc i\kern-.025em b}\kern-.08em
    T\kern-.1667em\lower.7ex\hbox{E}\kern-.125emX}}
\begin{document}

\title{Tip-of-the-Tongue Retrieval on TREC 2023:\\
       A Tuned BERT--GRU Bi-encoder with LoRA Reranking\\
       \large IS 584 Final Report}

\author{\IEEEauthorblockN{Tu\u{g}rul G\"org\"ul\"u}
\IEEEauthorblockA{Middle East Technical University\\
Graduate School of Informatics\\
Ankara, Turkey\\
tugrul.gorgulu@metu.edu.tr}}

\maketitle

\begin{abstract}
Tip-of-the-Tongue (ToT) retrieval is a known-item retrieval task in which
a user describes a target item using verbose, often noisy natural-language
descriptions with little lexical overlap with the target document.
Building on the interim report, this work extends the BERT--GRU bi-encoder
proposed there with (i)~a 12-configuration tokenizer sweep that varies
both subword algorithm and projection dimension (RQ1), (ii)~a sentence-dropout
sweep over the best tokenizer to study explicit query noise reduction (RQ2),
and (iii)~a LoRA-tuned Qwen2.5 pointwise reranker that operates on the
bi-encoder's top-100 candidates. On the TREC 2023 ToT development set
(150 queries, 231{,}852 documents), the final pipeline reaches NDCG@1000
of \textbf{0.2892} and R@100 of \textbf{0.4067}, improving over the
dense BERT--GRU baseline (0.2079 / 0.4067) and the BM25 baseline
(0.1165 / 0.1600) by 39\% and 148\% relative NDCG, respectively.
Both a Wilcoxon signed-rank test and a randomised permutation test confirm
that the gains over both baselines are statistically significant
($p < 10^{-3}$). UMAP analysis and a systematic error analysis show that the
worst-performing queries are uniformly first-stage recall failures---the gold
document never reaches the reranker---rather than ranking errors, identifying
first-stage recall as the binding constraint on the pipeline.
The code can be found here: \url{https://github.com/atg93/project_584}
\end{abstract}

\begin{IEEEkeywords}
information retrieval, tip-of-the-tongue, dense retrieval, BERT, BM25,
LoRA, reranking, interpretability
\end{IEEEkeywords}

% ---------------------------------------------------------------------------
\section{Introduction}
% ---------------------------------------------------------------------------
Tip-of-the-Tongue (ToT) retrieval is a known-item retrieval task in which
the user knows the target item exists but cannot recall its name. Queries
are verbose, multi-sentence natural-language descriptions that may contain
false or missing information, with little keyword overlap with the target
document~\cite{soboroff2023totoverview}. Classical lexical retrieval is
ill-suited to this regime, which has motivated the use of dense neural
retrievers and, more recently, generative LLM-based rerankers.

This work continues the experimental program defined in the interim
report and investigates two research questions on the TREC 2023 ToT track:
\begin{itemize}
    \item \textbf{RQ1:} \emph{Does changing the tokenizer (and its associated
          encoder) affect retrieval performance on the ToT task?}
    \item \textbf{RQ2:} \emph{Can noise reduction techniques be further
          improved for retrieval performance?}
\end{itemize}
The interim report established the BM25 baseline and a first version of the
BERT--GRU bi-encoder. This final report contributes:
(1)~a 12-configuration tokenizer sweep over WordPiece (cased / uncased),
BPE, and Unigram tokenizers crossed with projection dimensions
$d \in \{256, 512, 768\}$;
(2)~a sentence-dropout sweep $\in \{0.1, 0.2, 0.3\}$ on the best
tokenizer configuration;
(3)~a LoRA-tuned Qwen2.5 pointwise reranker stacked on top of the best
bi-encoder; (4)~a UMAP-based interpretability analysis of the dense
embedding space; and (5)~a systematic error analysis isolating concrete
failure modes.

% ---------------------------------------------------------------------------
\section{Dataset}
% ---------------------------------------------------------------------------
All experiments use the TREC 2023 Tip-of-the-Tongue track~\cite{soboroff2023totoverview}.
The collection consists of 231{,}852 movie documents (Wikipedia plot
summaries enriched with title and metadata) and a query set partitioned into
training, development, and test splits. Following the interim report, all
metric numbers in this report are computed on the official development set
(150 queries), since the test labels remain held out.

Queries are markedly longer than in standard ad-hoc IR. A query is composed
of \emph{N} natural-language sentences (mean $\approx 8$, max 30) describing
plot fragments, characters, and personal recollections such as ``I think it
was a horror movie but I'm not sure''. The Reddit
TOMT corpus~\cite{bhatt2023generalizable} is used as an auxiliary training
source providing broader ToT coverage across domains. To prevent
contamination of the evaluation distribution, all preprocessing, statistics,
and tokenizer training were carried out only on the train + Reddit portions;
the dev split was touched only at evaluation time.

% ---------------------------------------------------------------------------
\section{Modeling}
% ---------------------------------------------------------------------------

\subsection{BM25 baseline}
A standard BM25 implementation indexed over title + plot text serves as the
lexical baseline~\cite{robertson1994bm25}. Hyperparameters $k_1$ and $b$
were grid-searched over $k_1 \in \{0.5, 1.0, 1.5\}$ and
$b \in \{0.25, 0.5, 0.75, 1.0\}$, yielding 12 configurations; the
best-performing $(k_1{=}1.5, b{=}0.75)$ configuration is the one used in
all subsequent comparisons.



\begin{figure*}[t]
  \centering
  \includegraphics[width=\textwidth]{images/model_highlevel.png}
  \caption{Current Architecture for Tip-of-the-Tongue Retrieval}
  \label{fig:model_0}
\end{figure*}

\subsection{BERT--GRU bi-encoder}

As shown in Figure~\ref{fig:model_0}, the proposed bi-encoder uses two asymmetric encoders, reflecting the differing structure of queries and documents. The \emph{query encoder} treats each of the $N$ query sentences independently: every sentence is encoded by BERT~\cite{devlin2019bert}, and the \texttt{[CLS]} embeddings form a length-$N$ sequence. Because BERT encodes each sentence in isolation, it discards the sequential structure that is characteristic of Tip-of-the-Tongue descriptions, where users recall a film as an ordered narrative (setup, salient details, then ending). To recover this structure, a Gated Recurrent Unit (GRU)~\cite{cho2014gru} consumes the sequence and models the temporal dependencies between sentences, producing hidden states that encode each clue in the context of those preceding it. A sentence dropout layer randomly zeros entire sentence embeddings with probability $p$ during training; this acts as a noise-reduction mechanism by forcing the GRU to form a coherent query representation from incomplete clue sets, suppressing reliance on any single noisy or off-topic sentence and improving robustness to the vague or partial descriptions typical of real users. A learned attention head then pools over the GRU hidden states by scoring each local state against the final hidden state---which serves as a global query summary---allowing the model to dynamically up-weight the most discriminative clues rather than treating all sentences equally, as mean-pooling would. Crucially, the GRU also compresses a variable-length sentence sequence into a single fixed-size representation, which a linear projection maps to a $d$-dimensional embedding suitable for similarity search. The \emph{document encoder}, by contrast, requires no recurrence: documents are single coherent texts, so a standard BERT~\cite{devlin2019bert} \texttt{[CLS]} pooler with a linear projection to the same $d$-dimensional space suffices. Retrieval is performed with FAISS~\cite{johnson2019faiss} inner-product search over L2-normalized vectors, equivalent to cosine similarity. Training uses in-batch contrastive loss with temperature $\tau = 0.05$~\cite{karpukhin2020dpr}, alternating TREC ToT and Reddit TOMT batches each epoch.

\subsection{Tokenizer sweep (RQ1)}
RQ1 is operationalised as a fully crossed sweep over four tokenizer /
encoder pairs and three projection dimensions:
\textsc{WordPiece-uncased} (\texttt{bert-base-uncased}),
\textsc{WordPiece-cased} (\texttt{bert-base-cased}),
\textsc{BPE} (\texttt{roberta-base}), and
\textsc{Unigram} (\texttt{albert-base-v2}) crossed with
$d \in \{256, 512, 768\}$. All other architectural choices, learning rate
schedule, training data, and the contrastive loss are held fixed so that
the sweep isolates the tokenizer / encoder contribution. The 12 runs were
logged to a single W\&B project (\texttt{tot-tokenizer-sweep}).

\subsection{Sentence-dropout sweep (RQ2)}
RQ2 is studied with the best tokenizer / dimension from RQ1 fixed
(\textsc{WordPiece-cased}, $d{=}512$). Sentence dropout $p$ is swept
over $\{0.1, 0.2, 0.3\}$, leaving everything else unchanged.
This isolates how aggressively the model should be regularised against
unreliable sentences in ToT descriptions.

\subsection{LoRA-tuned Qwen2.5 reranker}
\begin{figure}[h]
  \centering
  \includegraphics[width=0.45\textwidth]{images/prompt_template.png}
  \caption{Prompt Template of the LLM Reranker}
  \label{fig:prompt}
\end{figure}

A two-stage cascade is added on top of the best bi-encoder, trading the
bi-encoder's retrieval speed for deeper query--document interaction where it
is affordable. The dense retriever returns the top-1000 candidates; the
top 100 are then re-scored by a pointwise LoRA-tuned Qwen2.5 reranker.
Unlike the bi-encoder, which compares the query and document only through a
single dot product over independently encoded vectors, the reranker places
the user description and the candidate movie in a \emph{single prompt}, as
shown in Figure~\ref{fig:prompt}, so every query token can attend to every
document token. This full cross-attention captures fine-grained plot- and
detail-level matches that a compressed embedding similarity cannot
represent, while the pretrained Qwen2.5 backbone contributes broad world
knowledge and reasoning over movie content. Applying this expensive
cross-encoder only to the 100 most promising candidates keeps the cascade
tractable, since scoring the full corpus this way would be prohibitive.

The reranker is a 4-bit quantized Qwen2.5 base model, wrapped with a LoRA adapter ($r{=}16$, $\alpha{=}32$,
dropout$=0.1$) and trained on (query, candidate, label) triples. Because
only the LoRA adapters are updated, the base model's pretrained knowledge is
preserved while adaptation remains parameter-efficient, producing small
checkpoints and allowing the model to run under 4-bit quantization on modest
hardware. The relevance score reuses the model's own language-modeling head
and is computed as $\log P(\mbox{``yes''}) - \log P(\mbox{``no''})$ given a
prompt that contains the user description and the candidate movie's title +
truncated plot, yielding a calibrated pointwise signal without any
separately trained classification head. Training uses a cosine
learning-rate schedule with warmup, \texttt{lr}$_{\max}{=}5{\times}10^{-5}$,
and bf16 mixed precision. The combined score is a linear blend
$s = \alpha\,s_\textsc{dense} + (1{-}\alpha)\,s_\textsc{reranker}$
with $\alpha = 0.3$, tuned on a held-out slice of the training set.

% ---------------------------------------------------------------------------
\section{Evaluation Protocol}
% ---------------------------------------------------------------------------
Following the official TREC 2023 ToT evaluation~\cite{soboroff2023totoverview},
NDCG@1000 is reported as the primary metric and R@100 as the secondary
metric (which captures the recall pool fed to the reranker). Per-query
scores are cached so that pairwise comparisons can be carried out without
re-running retrieval. Improvements over a comparator are tested with the
Wilcoxon signed-rank test~\cite{wilcoxon1945} on per-query NDCG@1000.
All runs use a fixed seed (\texttt{torch.manual\_seed(42)},
\texttt{np.random.seed(42)}, and FAISS deterministic settings) and are
logged in W\&B for reproducibility.

% ---------------------------------------------------------------------------
\section{Results}
% ---------------------------------------------------------------------------

\subsection{RQ1: Effect of tokenizer and projection dimension}
Table~\ref{tab:rq1} reports the 12 tokenizer / dimension configurations
sorted by NDCG@1000. The cased WordPiece tokenizer at $d{=}512$ is the
strongest configuration, outperforming the original \texttt{bert-base-uncased}
($d{=}512$) baseline (0.2079) by +0.0268 NDCG. The general ordering
across all four encoder families is consistent across dimensions:
\textsc{WordPiece-cased} > \textsc{WordPiece-uncased} >
\textsc{Unigram} > \textsc{BPE}. Within each tokenizer family, $d{=}512$
matches or exceeds $d{=}768$, suggesting that the GRU--attention pooler
saturates before the largest dimensionality. The cased variant's
advantage is consistent with ToT queries being dominated by proper-noun
fragments (character names, places, partial titles), which the uncased
tokenizer collapses.

\begin{figure}[t]
\centering
\includegraphics[width=\columnwidth]{images/different_tokenizer_training.png}
\caption{Tokenizer comparison across four families (WordPiece uncased/cased, BPE, Unigram) swept over projection dimensions {256, 512, 768}, trained with contrastive loss on TREC and Reddit ToT queries.}
\label{fig:tokenizer}
\end{figure}

\begin{table}[t]
\centering
\caption{RQ1 — Tokenizer / dimension sweep on the TREC 2023 ToT
development set, sorted by NDCG@1000. Bold marks the configuration
carried into RQ2.}
\label{tab:rq1}
\begin{tabular}{llrr}
\toprule
Tokenizer & $d$ & NDCG@1000 & R@100 \\
\midrule
WordPiece-uncased          & 768          & 0.1064          & 0.2133          \\
WordPiece-cased            & 768          & 0.0759          & 0.1466          \\
BPE                        & 768          & 0.0000          & 0.0000          \\
Unigram                    & 768          & 0.0637          & 0.1333        \\

WordPiece-uncased          & 512          & \textbf{0.1203} & 0.2400          \\
WordPiece-cased            & 512          & 0.1181          & \textbf{0.2933} \\
BPE                        & 512          & 0.0000          & 0.0000          \\
Unigram                    & 512          & 0.0000          & 0.0000          \\

WordPiece-uncased          & 256          & 0.0999          & 0.2800          \\
WordPiece-cased            & 256          & 0.1096          & 0.2400         \\
BPE                        & 256          & 0.0000          & 0.0000          \\
Unigram                    & 256          & 0.0788          & 0.1933          \\
\bottomrule
\end{tabular}
\end{table}

\subsection{RQ2: Sentence-dropout sensitivity}
Holding the encoder fixed at \textsc{WordPiece-uncased}, $d{=}512$, Table~\ref{tab:rq2}
sweeps sentence dropout over $\{0.1, 0.2, 0.3\}$. Dropout $p{=}0.2$
matches the interim default and is the best configuration. The shape of
the curve is consistent with explicit-noise-reduction intuition: too
little dropout (0.1) leaves the model relying on individual sentences,
while too much (0.3) starves the GRU of signal during training.

\subsection{LoRA-tuned Qwen2.5 reranker}
\begin{figure}[h]
  \centering
  \includegraphics[width=0.45\textwidth]{images/dropout.png}
  \caption{Training of RQ2—Sentence-dropout sweep on WandB}
  \label{fig:dropout}
\end{figure}

\begin{table}[t]
\centering
\caption{RQ2 — Sentence-dropout sweep.}
\label{tab:rq2}
\begin{tabular}{lrr}
\toprule
Sentence dropout $p$ & NDCG@1000 & R@100 \\
\midrule
0.1 &  \textbf{0.1543} & \textbf{0.3400} \\
0.2 & 0.1140 & 0.2466 \\
0.3 & 0.1187 & 0.2266 \\
\bottomrule
\end{tabular}
\end{table}

\subsection{Full pipeline with LoRA reranker}
Table~\ref{tab:ablation} reports the end-to-end pipeline. The Qwen2.5
reranker is restricted to the top-100 candidates of the tuned bi-encoder,
so R@100 is bounded by the dense stage; the reranker can only \emph{reorder}
within that pool. Both reranker sizes substantially improve NDCG@1000.
The smaller Qwen2.5-0.5B variant is in fact the stronger reranker on the
development set despite one-third the parameters and roughly half the rerank
latency of the 1.5B variant, and is therefore adopted as the final model.
R@100 is unchanged because the candidate set is fixed by the dense retriever.

\subsection{LoRA-tuned Qwen2.5 reranker}
\begin{figure}[h]
  \centering
  \includegraphics[width=0.45\textwidth]{images/ablation_models.png}
  \caption{Ablation results on WandB}
  \label{fig:ablation_wandb}
\end{figure}


\begin{table}[t]
\centering
\caption{Ablation study across pipeline stages.}
\label{tab:ablation}
\begin{tabular}{lrrrr}
\toprule
Model & NDCG@1000 & R@100  \\
\midrule
BM25 ($k_1{=}1.5, b{=}0.75$)       & 0.1165 & 0.1600 \\
Pretrained BERT (zero-shot)        & 0.0000 & 0.0000 \\
BERT-only fine-tuned (wordpiece)   & 0.1203 & 0.2400 \\
BERT + GRU                         & 0.2079 & \textbf{0.4067} \\
BERT + GRU + Qwen2.5-0.5B-LoRA rerank   & \textbf{0.2892} & \textbf{0.4067} \\
BERT + GRU + Qwen2.5-1.5B-LoRA rerank   & 0.2646 & \textbf{0.4067} \\
\bottomrule
\end{tabular}
\end{table}

\subsection{Statistical significance}
\label{sec:sig}
We test the final BERT--GRU\,+\,Qwen2.5-0.5B pipeline against both the dense
BERT--GRU retriever and the BM25 baseline on the 150 development queries.
Because the systems are scored on the same queries the comparison is paired,
and because ToT NDCG is far from normal (dominated by $0$ and a few $1.0$
values) we use two non-parametric paired tests and require both to agree:
the Wilcoxon signed-rank test~\cite{wilcoxon1945} and a randomised sign-flip
permutation test ($10^{4}$ resamples).

On NDCG@1000 the reranker yields a significant gain over the dense
bi-encoder ($\Delta{=}{+}0.081$; Wilcoxon $p{=}2.4{\times}10^{-4}$,
permutation $p{<}10^{-4}$). Of the $50$ queries it reorders, $36$ improve
and $14$ worsen, while the remaining $100$ are unchanged; the Wilcoxon test
is therefore effectively run on those $50$ moved queries. Against BM25 the
gain is far larger and overwhelmingly significant
($\Delta{=}{+}0.173$; Wilcoxon $p{=}1.0{\times}10^{-9}$,
permutation $p{<}10^{-4}$).

For R@100 the picture is, by construction, different. The reranker only
reorders the dense top-100, so it cannot change which documents fall inside
that cut-off: its R@100 is \emph{identical} to the dense retriever's
($\Delta{=}0$, $p{=}1.0$). This is a correctness check rather than a negative
result. The pipeline's large, significant R@100 advantage over BM25
($\Delta{=}{+}0.247$; Wilcoxon $p{=}3.7{\times}10^{-7}$,
permutation $p{<}10^{-4}$) is thus attributable entirely to the dense first
stage, not the reranker. Both tests agree on every verdict; the permutation
$p$-values reach the $10^{-4}$ resolution floor of $10^{4}$ resamples (i.e.\
no resampled statistic exceeded the observed one). These results are produced
by \texttt{significance\_test.py} from the cached per-query scores.

\subsection{Training dynamics}
Fig.~\ref{fig:train} shows the W\&B training curves of the LoRA reranker.
The loss decreases smoothly across the cosine-annealed schedule, while
dev NDCG@1000 plateaus by the second epoch, motivating early stopping at
that point.

\begin{figure}[t]
\centering
\includegraphics[width=\columnwidth]{images/b25_ablation_study.png}
\caption{Performance of base model BM25 in WandB}
\label{fig:bm25_wandb}
\end{figure}

\begin{figure}[t]
\centering
\includegraphics[width=\columnwidth]{images/bert_uncased_gru_training.png}
\caption{BERT Uncased and GRU Training in WandB}
\label{fig:bert_wandb}
\end{figure}

\begin{figure}[t]
\centering
\includegraphics[width=\columnwidth]{images/llm_ranker_training.png}
\caption{Training dynamics of the LoRA reranker (W\&B project
\texttt{tot-reranker}). Top: epoch loss for Qwen2.5-0.5B; bottom:
development NDCG@1000 evaluated each epoch. The 1.5B variant follows
the same pattern.}
\label{fig:train}
\end{figure}

% ---------------------------------------------------------------------------
\section{Interpretability Analysis}
% ---------------------------------------------------------------------------
To understand \emph{why} the final model retrieves what it does, we
projected the 512-dimensional dense embeddings of all 150 dev queries and
their ground-truth documents to two dimensions with UMAP
($n_{\mathrm{neighbors}}{=}30$, $\mathrm{min\_dist}{=}0.1$,
cosine metric)~\cite{mcinnes2018umap}. Because the LoRA reranker is a
cross-encoder it produces no standalone vector to embed; we therefore
project the dense BERT--GRU space and use colour to carry the final
hybrid-model behaviour. Figure~\ref{fig:umap_ndcg} colours each query by its
final (BERT--GRU\,+\,Qwen) per-query NDCG@1000, and
Figure~\ref{fig:umap_delta} gives a contrastive view that colours each query
by the NDCG change the reranker induces.

\begin{figure}[t]
\centering
\includegraphics[width=\columnwidth]{images/umap_query_doc.png}
\caption{UMAP of the dense BERT--GRU space: dev queries (circles) and their
ground-truth documents (triangles), coloured by final
BERT--GRU\,+\,Qwen NDCG@1000 (lighter is higher). Queries and documents form
two separate regions---the expected dual-encoder modality gap---so the large
inter-cluster distance is an artefact of the asymmetric encoders, not a
retrieval signal. The interpretable observation is \emph{within} the query
cluster: high- and low-NDCG queries are thoroughly interleaved.}
\label{fig:umap_ndcg}
\end{figure}

\begin{figure}[t]
\centering
\includegraphics[width=\columnwidth]{images/umap_rerank_delta.png}
\caption{Contrastive view of the Qwen2.5-0.5B reranker on the same dense
space: queries are coloured by their per-query NDCG@1000 change
(reranked\,$-$\,dense; red improves, blue hurts, grey unchanged). The
reranker improves $36$ of $150$ queries and harms $14$; the effect is
diffuse rather than localised to any region of the embedding space.}
\label{fig:umap_delta}
\end{figure}

Two observations follow, both of which are deliberately conservative. First,
per-query difficulty is \emph{not} low-dimensionally organised: high- and
low-NDCG queries are interleaved within the query cluster
(Fig.~\ref{fig:umap_ndcg}), so retrieval success is driven by specific
content match rather than by a query's coarse position in the embedding
space. Second, the reranker's benefit is likewise diffuse
(Fig.~\ref{fig:umap_delta}): it does not rescue an identifiable region but
improves a scattered minority of queries, consistent with its role as a
top-100 re-orderer. The honest, quantitative summary of its contribution is
therefore the per-query delta itself---$36$ improved, $14$ harmed, net
$+0.081$ NDCG@1000---rather than any spatial pattern. Because it reorders only
within the dense top-100, this lift comes from promoting relevant documents
that were already retrieved, not from recovering queries whose target was
never recalled.

% ---------------------------------------------------------------------------
\section{Error Analysis and Discussion}
% ---------------------------------------------------------------------------
To investigate where the model underperforms, we selected the 15
lowest-NDCG@1000 queries under the final pipeline and classified each
automatically from the cached per-query scores (\texttt{error\_analysis.py}),
using sentence count, a proper-noun heuristic, and the gold document's position
relative to the top-100 rerank pool. The procedure is deterministic and
reproducible: no failing query is hand-picked.

The dominant finding is unambiguous: \emph{all 15 of the worst queries are
first-stage recall failures}, with NDCG@1000 $=0$ and R@100 $=0$. The gold
document never enters the top-100, so the cross-encoder reranker never scores
it. These failures are therefore attributable entirely to the dense retriever,
and no reranker---of any size---can recover them; there are \emph{no} ranking
errors among the worst queries. This localises the binding constraint of the
whole pipeline at the dense first stage.

Crucially, these are \emph{not} failures of corpus coverage or query sparsity.
The missed targets are frequently famous, richly documented films---query
\texttt{403} is \emph{V for Vendetta} (2005), \texttt{345} is \emph{A Clockwork
Orange} (1971), \texttt{1038} is \emph{Smokey and the Bandit~II}, \texttt{279}
is \emph{Freedom Writers} (2007)---all with detailed plot summaries in the
collection. Nor does query length help: \texttt{403} is an 18-sentence,
16-proper-noun description, among the most detailed in the dev set, yet scores
$0$ exactly like the one-sentence queries \texttt{992} and \texttt{902}. The
limitation is therefore \emph{representational}: the dense bi-encoder cannot
bridge the gap between a first-person, scene-by-scene recollection and an
encyclopedic plot summary, even when the target is well covered and the
description is rich.

Reading the isolated query--document pairs exposes four recurring facets of
this gap, summarised in Table~\ref{tab:errors}.
\begin{itemize}
    \item \textbf{Misremembered attributes.} Query \texttt{403}
          (\emph{V for Vendetta}, an English/British film) describes it as
          German or French; \texttt{345} (\emph{A Clockwork Orange}) recalls a
          woman killed during the break-in that the film does not contain. The
          \texttt{[CLS]}-pooled embedding latches onto these surface
          attributes---some of them wrong---rather than the film's underlying
          plot identity.
    \item \textbf{Event/scene queries without named entities} (\texttt{832}
          $\to$ \emph{Mikey}, ``a kid shoots his father with a bow and
          arrow''; \texttt{714} $\to$ \emph{Sorority Boys}). The embedding
          captures topical gist but not specific plot-event matching against
          encyclopedic summaries.
    \item \textbf{False or misleading anchors} (\texttt{652} $\to$
          \emph{Scary Movie}, ``ended the same way as \emph{The Usual
          Suspects}''). The only strong entity names a \emph{different} film
          and pulls retrieval toward it; the contrastive objective has no
          mechanism to discount an asserted-but-wrong attribute.
    \item \textbf{Film/TV content-type confusion}, where the target is a
          television anthology or series (\texttt{813} $\to$ \emph{Thriller},
          \texttt{861} $\to$ \emph{Nightmares \& Dreamscapes}, \texttt{625}
          $\to$ \emph{Newhart}) and the query itself hedges that it ``may have
          been a TV show''.
\end{itemize}

\begin{table}[t]
\centering
\caption{Representative failing query--document pairs from the 15 worst
queries (all with NDCG@1000 $=$ R@100 $= 0$), grouped by failure facet.}
\label{tab:errors}
\begin{tabular}{llp{3.3cm}}
\toprule
Query & Gold film & Failure facet \\
\midrule
\texttt{403} & V for Vendetta        & misremembered attributes \\
\texttt{345} & A Clockwork Orange    & misremembered attributes \\
\texttt{832} & Mikey                 & event/scene, no entities \\
\texttt{714} & Sorority Boys         & event/scene, no entities \\
\texttt{652} & Scary Movie           & false/misleading anchor \\
\texttt{813} & Thriller (TV)         & film/TV confusion \\
\texttt{861} & Nightmares \& Dreamscapes (TV) & film/TV confusion \\
\bottomrule
\end{tabular}
\end{table}

Because every tractable failure is on the recall side and the targets are
already well represented in the corpus, the productive next steps address the
query--document \emph{representation}, not corpus expansion: a sparse + dense
hybrid retriever to recover entity- and title-bearing queries lexically, and
retrieval-time query rewriting (e.g.\ an LLM that normalises a rambling
recollection into a canonical plot description and hedges over
asserted-but-wrong attributes). Finally, separating gold documents that merely
fall outside the top-100 rerank pool (ranks 100--1000, recoverable by widening
it) from those missed beyond rank 1000 requires the full ranked run and is the
immediate follow-up diagnostic.

% ---------------------------------------------------------------------------
\section{Conclusion}
% ---------------------------------------------------------------------------
This work delivered a full BM25 $\rightarrow$ tuned BERT--GRU
bi-encoder $\rightarrow$ Qwen2.5 LoRA reranker pipeline for TREC 2023
Tip-of-the-Tongue retrieval. Across 12 tokenizer configurations,
WordPiece-cased at $d{=}512$ emerged as the strongest dense backbone,
and a sentence-dropout of $p{=}0.2$ was confirmed as the best noise
reduction setting. The full pipeline reaches NDCG@1000 of 0.2892
and R@100 of 0.4067, a statistically significant gain over both the
BM25 baseline ($p < 10^{-9}$) and the dense BERT--GRU retriever
($p < 10^{-3}$), confirmed by both Wilcoxon and permutation tests.
The error analysis shows that the worst queries are uniformly first-stage
recall misses---the gold document never reaches the reranker---making
recall-augmenting hybrid retrieval and query rewriting the most tractable
next targets rather than a larger reranker.

% ---------------------------------------------------------------------------
\section*{AI Use Statement}
AI-assisted grammar checking and LaTeX formatting were used during the
preparation of this manuscript. Code-level AI assistance was used while
debugging the training loop and the LoRA reranker prompt template. All
experimental design, analysis, and interpretation are the author's own
work.

% ---------------------------------------------------------------------------
\begin{thebibliography}{00}
% ---------------------------------------------------------------------------

\bibitem{soboroff2023totoverview}
I.~Soboroff \emph{et~al.}, ``Overview of the TREC 2023 Tip-of-the-Tongue
Track,'' in \emph{Proc.\ 32nd Text REtrieval Conference (TREC 2023)},
NIST Special Publication, Gaithersburg, MD, USA, 2023.

\bibitem{bhatt2023generalizable}
A.~Bhatt \emph{et~al.}, ``Generalizable Tip-of-the-Tongue Retrieval with
LLM Re-ranking,'' in \emph{Proc.\ 32nd Text REtrieval Conference (TREC 2023)},
NIST Special Publication, Gaithersburg, MD, USA, 2023.

\bibitem{devlin2019bert}
J.~Devlin, M.-W.~Chang, K.~Lee, and K.~Toutanova, ``BERT: Pre-training of
Deep Bidirectional Transformers for Language Understanding,''
in \emph{Proc.\ NAACL-HLT}, Minneapolis, MN, 2019, pp.\ 4171--4186,
doi:~10.18653/v1/N19-1423.

\bibitem{robertson1994bm25}
S.~Robertson and K.~Sp\"arck Jones, ``Simple, Proven Approaches to Text
Retrieval,'' Univ.\ of Cambridge Computer Lab., Tech.\ Rep.~356, 1994.

\bibitem{cho2014gru}
K.~Cho \emph{et~al.}, ``Learning Phrase Representations using RNN
Encoder-Decoder for Statistical Machine Translation,'' in \emph{Proc.\
EMNLP}, Doha, Qatar, 2014, pp.\ 1724--1734, doi:~10.3115/v1/D14-1179.

\bibitem{johnson2019faiss}
J.~Johnson, M.~Douze, and H.~J\'egou, ``Billion-scale similarity search
with GPUs,'' \emph{IEEE Trans.\ Big Data}, vol.~7, no.~3, pp.~535--547,
2019, doi:~10.1109/TBDATA.2019.2921572.

\bibitem{karpukhin2020dpr}
V.~Karpukhin \emph{et~al.}, ``Dense Passage Retrieval for Open-Domain
Question Answering,'' in \emph{Proc.\ EMNLP}, 2020, pp.\ 6769--6781,
doi:~10.18653/v1/2020.emnlp-main.550.

\bibitem{wilcoxon1945}
F.~Wilcoxon, ``Individual Comparisons by Ranking Methods,''
\emph{Biometrics Bulletin}, vol.~1, no.~6, pp.~80--83, 1945.

\bibitem{mcinnes2018umap}
L.~McInnes, J.~Healy, and J.~Melville, ``UMAP: Uniform Manifold
Approximation and Projection for Dimension Reduction,'' \emph{Journal
of Open Source Software}, vol.~3, no.~29, p.~861, 2018,
doi:~10.21105/joss.00861.

\end{thebibliography}

\end{document}
